\begin{longtable}{@{}P{4.3cm}rrrP{6.0cm}@{}}
\toprule
\textbf{Parámetro} & \textbf{Mín.} & \textbf{Moda} & \textbf{Máx.} & \textbf{Fuente o justificación} \\
\midrule
\endhead
IDM sin SGC -- Viguetas & 0,0~\% & 1,0~\% & 3,0~\% & Supuesto del autor: viguetas fabricadas a medida (Plano 03); pérdidas por rotura en manipulación \\
IDM sin SGC -- Plastoformo & 3,0~\% & 6,0~\% & 9,0~\% & provisional \\
IDM sin SGC -- Hormigón & 5,0~\% & 10,0~\% & 22,0~\% & Galarza Meza (2011, pp. 53, 73): 5 % típico y 22 % sin control; Ruiz Muñoz y Rojas Bonilla (2018, pp. 33, 40): 14–27 % semanal real sin control; Maciel et al. (2016, p. 106): 1,5–13 %. Moda: valor intermedio para obra sin control volumétrico \\
IDM sin SGC -- Acero & 4,0~\% & 10,0~\% & 17,1~\% & Galarza Meza (2011, p. 75): 7 % presupuestado; Abbott y Rodríguez (2008, p. 583): hasta 10 % cortado en obra; Chino Quispe y Zarate Chilqquetuma (2024): 14,58–17,09 % con método convencional \\
Eficacia del SGC -- Viguetas & 10,0~\% & 44,0~\% & 70,0~\% & Igual que materiales de obra (Galarza Meza, 2011, p. 73) \\
Eficacia del SGC -- Plastoformo & 10,0~\% & 44,0~\% & 70,0~\% & Igual que materiales de obra (Galarza Meza, 2011, p. 73) \\
Eficacia del SGC -- Hormigón & 10,0~\% & 44,0~\% & 70,0~\% & Galarza Meza (2011, p. 73): de 5 % a 2,8 % con pedidos limitados (44 %); Dzul y Gracia (2009, p. 46): 10 % \\
Eficacia del SGC -- Acero & 10,0~\% & 42,0~\% & 80,0~\% & Galarza Meza (2011, pp. 75-76): de 7 % a 4,04 % (42 %) y 3,75 % (46 %); Abbott y Rodríguez (2008, p. 583): < 1 % precortado \\
Fracción de retrabajo sin SGC ($\rho$) & 0,5~\% & 2,3~\% & 6,5~\% & Khadim et al. (2023, p. 5421): fallas visibles 2,34 %; Gudiel Cairampoma y Solis Puente (2020, p. 1): no calidad 2,7–5,1 % de la partida de concreto armado; Álvarez Quintero (2014, p. 12): 6,5 % sin SGC \\
Eficacia del SGC en retrabajos ($e_R$) & 10,0~\% & 49,0~\% & 89,0~\% & Dzul y Gracia (2009, p. 46): 10 %; Palanisamy y Palanichamy (2025): 49 %; Álvarez Quintero (2014, p. 12): de 6,5 % a 0,72 % (89 %) \\
Encargado de calidad a tiempo parcial (3 meses) (Bs) & 4.000 & 6.750 & 13.500 & 50 % de 3 meses de un jefe de calidad con sueldo anual de Bs 54.000 (Peña Callisaya, 2024, p. 171); máximo: 100 % \\
Ensayos de hormigón (probetas y asentamiento) (Bs) & 720 & 1.200 & 1.920 & 24 probetas según ACI 318-25 (26.12.2.1) a Bs 30–80 por probeta (supuesto) \\
Capacitación del personal (3 charlas) (Bs) & 1.000 & 1.700 & 5.874 & 12 trabajadores × 4,5 h con tarifas de Cocarico Mamani (2022); máximo: Peña Callisaya (2024, p. 171) \\
Control de almacén y conciliación de cantidades (Bs) & 500 & 820 & 1.500 & 0,5 h/día durante 60 días con tarifa de ayudante con cargas (Cocarico Mamani, 2022) \\
Formatos, registros e impresión (Bs) & 200 & 400 & 1.000 & Supuesto del autor \\
\bottomrule
\end{longtable}
